\documentclass[standalone]{beamer}

\begin{document}
\section{組合}

\begin{frame}{\btitle{目標}}
  \begin{problem}[經典問題]
    一個項鍊由 $N$ 個寶石串成，有 $M$ 種不同的寶石，並且因為項鍊是環形的，所以旋轉相同視為相同的。
    請問有多少種不同的項鍊？
  \end{problem}
  \pause

  \begin{enumerate}
    \item 有很多\emph{物體}構成一個集合 $X$。
    \item 有一些\emph{旋轉} $G$。
    \item 旋轉相同視為相同，也就是旋轉會把物件劃分成許多等價類。
  \end{enumerate}
\end{frame}

\newcommand{\chain}[6][]{
  \begin{scope}[#1 shift={#2}, w/.style={fill=white}, b/.style={fill=black}, nd/.style={draw, circle}]
    \draw (0, 0) circle (0.3);
    \node[nd] at (0:0.3) [#3] {};
    \node[nd] at (90:0.3) [#4] {};
    \node[nd] at (180:0.3) [#5] {};
    \node[nd] at (270:0.3) [#6] {};
  \end{scope}
}
\begin{frame}{\btitle{目標}}
  \begin{columns}
    \begin{column}{0.6\textwidth}
      \begin{figure}
        \begin{tikzpicture}[x=1.3cm, y=-1.3cm]
          \chain{(0, 0)}{w}{w}{w}{w}
          \chain{(3, 0)}{b}{b}{b}{b}

          \chain{(0, 1)}{w}{w}{w}{b}
          \chain{(1, 1)}{w}{w}{b}{w}
          \chain{(2, 1)}{w}{b}{w}{w}
          \chain{(3, 1)}{b}{w}{w}{w}

          \chain{(0, 2)}{b}{w}{b}{w}
          \chain{(1, 2)}{w}{b}{w}{b}

          \chain{(0, 3)}{b}{b}{w}{w}
          \chain{(1, 3)}{w}{w}{b}{b}

          \chain{(0, 4)}{b}{b}{b}{w}
          \chain{(1, 4)}{b}{b}{w}{b}
          \chain{(2, 4)}{b}{w}{b}{b}
          \chain{(3, 4)}{w}{b}{b}{b}
        \end{tikzpicture}
        \caption{$M = 2$, $N = 2$ 的例子}
      \end{figure}
    \end{column}
    \vrule{}
    \begin{column}{0.4\textwidth}
      \pause
      \begin{itemize}[<+->]
        \itemsep=1ex
        \item 物體有 $\abs{X} = 2^4 = 16$ 種。
        \item 令 $a$ 為旋轉 $90^\circ$，則
          $G = \{1, a, a^2, a^3\}$。
        \item 最後答案是 $5$。
      \end{itemize}
    \end{column}
  \end{columns}
\end{frame}

\begin{frame}
  \begin{definition}
    \begin{enumerate}
      \item $G_x \defeq \{ g \in G : g x = x \}$，也就是固定 $x$ 下，所有不會動到 $x$ 的作用。
      \item $X^g \defeq \{ x \in X : g x = x \}$，也就是固定一個作用 $g$ 下的\emph{不動點}。
      \item $G x \defeq \{ gx : g \in G \}$，也被稱作是 $x$ 在 $G$ 下的\emph{軌道}。
      \item $X / G \defeq \{ Gx : x \in X \}$，也就是 $X$ 在 $G$ 下所有的軌道。
    \end{enumerate}
  \end{definition}
\end{frame}

\begin{frame}{\btitle{目標}}
  \begin{columns}
    \begin{column}{0.6\textwidth}
      \begin{figure}
        \begin{tikzpicture}[x=1.3cm, y=-1.3cm, remember picture]
          \chain{(0, 0)}{w}{w}{w}{w}
          \chain{(3, 0)}{b}{b}{b}{b}

          \chain{(0, 1)}{w}{w}{w}{b}
          \chain{(1, 1)}{w}{w}{b}{w}
          \chain{(2, 1)}{w}{b}{w}{w}
          \chain{(3, 1)}{b}{w}{w}{w}

          \chain{(0, 2)}{b}{w}{b}{w}
          \chain{(1, 2)}{w}{b}{w}{b}
          \node (burn_x) at (1.4, 2) {};

          \chain{(0, 3)}{b}{b}{w}{w}
          \chain{(1, 3)}{w}{w}{b}{b}

          \chain{(0, 4)}{b}{b}{b}{w}
          \chain{(1, 4)}{b}{b}{w}{b}
          \chain{(2, 4)}{b}{w}{b}{b}
          \chain{(3, 4)}{w}{b}{b}{b}
        \end{tikzpicture}
      \end{figure}
    \end{column}
    \vrule{}
    \begin{column}{0.4\textwidth}
      \pause
      令 \tikzoverlay{burn_t}{$x$} 為箭頭所指的物體，
      $a$ 為旋轉 $90^\circ$。
      \begin{itemize}[<+->]
        \itemsep=1ex
        \item $G_x = \alt<+->{\{1, a^2\}}{\ ?}$
        \item $X^{a^2} = \ ?$
        \item $G_x = \ ?$
        \item $x^G = \ ?$
        \item $X/G = \ ?$
      \end{itemize}
      \begin{tikzpicture}[overlay, remember picture]
        \draw[-latex] ($(burn_t)+(0, 0.2)$) ..controls ($(burn_t) + (0, 1)$) and ($(burn_t) + (-1.0, 0.7)$) .. ($(burn_t) + (-1.5, 0)$)
         ..controls ($(burn_t) + (-2.0, -0.7)$) and ($(burn_x) + (3.0, 0)$) .. (burn_x);
      \end{tikzpicture}
    \end{column}
  \end{columns}
\end{frame}

\begin{frame}{\btitle{目標}}
\begin{theorem}[Burnside lemma]
  \[ \abs{X / G} = \frac{1}{\abs{G}} \sum_{g \in G} \abs{X^g} \]
\end{theorem} \pause
也就是說要算把旋轉相同視為相同的個數：\pause \disskip
\begin{enumerate}[<+->]
  \item 對於每個旋轉 $g$，計算在這個旋轉下的不動點的個數 $\abs{X^g}$。
  \item 把這些數字加起來除以 $\abs{G}$ 就是答案。
\end{enumerate}
\end{frame}

\begin{frame}

  \begin{columns}
    \begin{column}{0.6\textwidth}
      \begin{figure}
        \begin{tikzpicture}[x=1.3cm, y=-1.3cm, remember picture]
          \chain{(0, 0)}{w}{w}{w}{w}
          \fill[blue, opacity=0.2] (0, 0) circle (1);
          \chain{(3, 0)}{b}{b}{b}{b}

          \chain{(0, 1)}{w}{w}{w}{b}
          \chain{(1, 1)}{w}{w}{b}{w}
          \chain{(2, 1)}{w}{b}{w}{w}
          \chain{(3, 1)}{b}{w}{w}{w}

          \chain{(0, 2)}{b}{w}{b}{w}
          \chain{(1, 2)}{w}{b}{w}{b}

          \chain{(0, 3)}{b}{b}{w}{w}
          \chain{(1, 3)}{w}{w}{b}{b}

          \chain{(0, 4)}{b}{b}{b}{w}
          \chain{(1, 4)}{b}{b}{w}{b}
          \chain{(2, 4)}{b}{w}{b}{b}
          \chain{(3, 4)}{w}{b}{b}{b}
        \end{tikzpicture}
        \caption{$M = 2$, $N = 2$ 的例子}
      \end{figure}
    \end{column}
    \vrule{}
    \begin{column}{0.4\textwidth}
      \begin{enumerate}[<+->]
        \item $\abs{X^{1}} = \onslide<+->{16} $
      \end{enumerate}
    \end{column}
  \end{columns}
\end{frame}

\end{document}
